\section{Heartbeats}

En un sistema distribuido un nodo puede fallar sin que el resto lo note
de inmediato. Hay que distinguir un nodo caído de uno lento. Usamos
heartbeats y timeouts: cada nodo envía periódicamente un ``estoy vivo''
y, si el monitor no lo recibe a tiempo, lo declara fallido. El módulo
\texttt{heartbeat} simula eso. Se ejecuta con
\texttt{cargo run -- parte2}. Con \texttt{parte2-lider} se fuerza la
caída del líder para disparar la elección de la sección siguiente.

El clúster tiene tres nodos. Cada uno envía un \texttt{Heartbeat} cada
500~ms por un canal \texttt{mpsc} hacia un monitor. Por cada nodo hay un
vigilante: un hilo que ejecuta \texttt{vigilar} y hace
\texttt{recv\_timeout} sobre el receptor correspondiente. El tiempo de
espera depende del estado actual. Desde \texttt{Vivo} espera el umbral
de sospecha. Si no llega nada, pasa a \texttt{Sospechoso} y emite
\texttt{Evento::Sospechoso}. Desde \texttt{Sospechoso} espera la
diferencia entre el umbral de caída y el de sospecha; si el silencio
continúa, pasa a \texttt{Caido}, emite \texttt{Evento::Caido} y termina.
Si llega un latido estando sospechoso, vuelve a \texttt{Vivo} y emite
\texttt{Evento::Recuperado}, porque el monitor se había equivocado.

Los umbrales son 1~s para sospechoso y 2~s para caído, con latidos cada
500~ms. \texttt{Config::rapida()} usa los mismos valores divididos por
diez (50, 100 y 200~ms) para que las pruebas no tarden de más.

\subsection{Caída del nodo}

La caída simulada no cierra el canal. El nodo deja de enviar porque un
\texttt{AtomicBool} llamado \texttt{caido} pasa a verdadero, pero el
\texttt{Sender} permanece vivo. Si el nodo terminara de forma limpia,
\texttt{recv\_timeout} devolvería
\texttt{RecvTimeoutError::Disconnected} y el vigilante saldría del bucle
sin declarar falla: sería un cierre ordenado. Si el nodo se cuelga o la
red se parte, el canal queda abierto y no llega nada; el silencio es lo
único que el monitor ve. El flag modela un servidor detenido, no un
proceso que terminó bien.

\texttt{simular} inyecta la falla cuando transcurre \texttt{caer\_tras},
deja correr el tiempo hasta que el vigilante declare la caída y corta en
el primer \texttt{Evento::Caido} de la víctima. Los otros dos nodos
siguen enviando latidos. La prueba verifica que nunca se los declare
caídos.

\subsection{Detector imperfecto}

Un detector basado en timeouts no distingue un nodo caído de uno lento.
En un sistema asíncrono no existe un detector de fallas perfecto, porque
el silencio puede deberse a una caída o a una red lenta. Con un umbral
corto aumentan los falsos positivos. Con un umbral largo se demora más
la reacción ante una falla real. El monitor solo registra que el latido no llegó a tiempo.

\subsection{Pruebas}

Con tres nodos, al detener el envío del nodo~1 aparece primero
\texttt{Sospechoso(1)} y después \texttt{Caido(1)}, sin declarar caídos
a los demás. Un vigilante que recibe cinco latidos a tiempo no emite
ningún evento. Si el nodo se silencia 150~ms (entre el umbral de sospecha
y el de caída) y luego envía un latido, la secuencia es
\texttt{Sospechoso(0)} y \texttt{Recuperado(0)}. Las tres pruebas de
\texttt{cargo test --lib heartbeat} pasan.
